% LaTeX report template
%

\documentclass{article}
\usepackage[a4paper]{geometry}
\usepackage{graphicx}
\usepackage[english]{babel}
\usepackage[latin1]{inputenc}
\usepackage{caption}
%
\usepackage{hyperref}
\hypersetup{
	colorlinks,
	citecolor=blue,
	filecolor=black,
	linkcolor=blue,
	urlcolor=blue
}
\usepackage{amsfonts,graphicx, tikz, amssymb, tikz-cd}
\usepackage{flafter}
\usepackage{dcolumn}
\usepackage{booktabs}
\usetikzlibrary{positioning,shapes,arrows,shapes.geometric}
\usetikzlibrary{bayesnet}
\newcolumntype{M}[1]{D{.}{.}{1.#1}}
\setlength{\textwidth}{6.25in}          % normal width of text on page
\setlength{\oddsidemargin}{0in}   % left margin less 1 inch, odd #'ed pgs
\setlength{\evensidemargin}{.25in}
\usetikzlibrary{backgrounds,automata}
\usetikzlibrary{decorations.pathreplacing}
\usetikzlibrary{calc}
\tikzstyle{graphnode}=[circle,draw, fill=white, minimum size=2em]

\usepackage{algorithm}
\usepackage[noend]{algpseudocode}
\newcommand{\mopbot}{MopBot}
\usepackage{booktabs}
\newcommand{\myindent}[1]{
	\newline\makebox[#1cm]{}
}
\newcommand{\lmodels}{\rotatebox[origin=c]{180}{$\models$}}
\usepackage{environ}
\usepackage[fleqn]{amsmath}

\newtheorem{theorem}{Theorem}
\newtheorem{example}{Example}
\newtheorem{lemma}[theorem]{Lemma}
\newtheorem{proposition}[theorem]{Proposition}
\newtheorem{claim}[theorem]{Claim}
\newtheorem{corollary}[theorem]{Corollary}
\newtheorem{definition}[theorem]{Definition}
\newenvironment{proof}{{\bf Proof:}}{\hfill\rule{2mm}{2mm}}
\usepackage[stable]{footmisc}

\DeclareMathOperator*{\argmin}{\arg\!\min}
\DeclareMathOperator*{\argmax}{\arg\!\max}


\newcommand{\ap}[1]{\ensuremath{\mathsf{AP}{#1}}}
\tikzset{%
	zeroarrow/.style = {-stealth,dashed},
	onearrow/.style = {-stealth,solid},
	c/.style = {circle,draw,solid,minimum width=2em,
		minimum height=2em},
	r/.style = {rectangle,draw,solid,minimum width=2em,
		minimum height=2em}
}
\begin{document}
\title{Lecture 2: Satisfiability}

\section{Proposition Logic: Satisfiability}

\subsection{Relevance Lemma}
Now, let us consider a formula $\varphi:= (p_1 \lor p_2)$, and the following truth assignments:
\begin{enumerate}
	\item $\tau_1: \{p_1 \mapsto 1, p_2 \mapsto 1, p_3 \mapsto 0, p_4 \mapsto 1 \}$.
	\item $\tau_2: \{p_1 \mapsto 1, p_2 \mapsto 1, p_3 \mapsto 1, p_4 \mapsto 0 \}$.
	\item $\tau_3: \{p_1 \mapsto 1, p_2 \mapsto 1, \ldots\}$.
\end{enumerate}

Clearly, we see that $\varphi(\tau_1)$ and $\varphi(\tau_2)$ is $1$, but what about $\tau_3$? Observe that irrespective of valuation of $p_3$, and $p_4$, we can say that $\varphi(\tau_3) = 1$. That is, we want to capture the notation of only relevant propositions instead of considering the entire PROP.

\begin{definition}
	The set of Atomic Propositions \ap{(\varphi)} is defined as follows:
	\begin{itemize}
		\item \ap{(p)} = \{p\}
		\item \ap{(\lnot \alpha)} = \ap{(\alpha)}
		\item \ap{(\alpha \odot \beta)} = \ap{(\alpha)} $\cup$ \ap{(\beta)},  where $\odot$ is a binary connective.
	\end{itemize}
\end{definition}

\begin{lemma}[Relevance lemma]
	Let $\varphi$ be the formula and $\tau_1, \tau_2 \in 2^{\text{PROP}}$, $$\tau_1 \cap \ap{(\varphi)} = \tau_2 \cap \ap{(\varphi)} \implies \varphi(\tau_1) = \varphi(\tau_2)$$
\end{lemma}

Relevance lemma tells us that we need to \textit{only} care about $2^{\ap{(\varphi)}}$.

\subsection{Set of all Models}

\begin{definition}
	Let $\varphi$ be a formula and the set of all models of $\varphi$ is defined as:
	$$\text{models}(\varphi) = \{\tau | \varphi(\tau) = 1 \}$$

\end{definition}

The above definition of models does not give recipe for computation of models($\varphi$).
\begin{enumerate}
	\item models(p) = $2^{\text{PROP}} \setminus {\lnot p}$.
	\item models($\lnot \alpha$) = $2^{\text{PROP}} \setminus \text{models}(\alpha)$.
	\item models($\alpha \lor \beta$) = models($\alpha$) $\cup$ models($\beta$).
	\item models($\alpha \land \beta$) = models($\alpha$) $\cap$ models($\beta$).
\end{enumerate}
Example: $\varphi := (p \lor q) \; \text{models}(\varphi) = \{ \{p\},  \{q\}, \{p,q\} \}$.
We can view as for every formula $\varphi$, there is a unique models($\varphi$) $\subseteq 2^{\text{PROP}}$.
\begin{definition}
	Two formulas $\varphi$, and $\psi$ is considered to be logically equivalent to each other if and only if they both have the same set of models, i.e., $\varphi \equiv \psi$ if and only if models($\varphi$) = models($\psi$).
\end{definition}
Example, let $\varphi := (p)$, and $\psi := (\lnot(\lnot p))$. Observe that models($\varphi$) = models($\psi$) = $\{\{p\}\}$, therefore $(p) \equiv (\lnot (\lnot p))$.


\section{Resolution}

\subsection{Conjunctive Normal Form}

The Conjunctive Normal Form (CNF) of a formula $C$ is represented as:
\begin{align*}
	C = C_1 \land C_2 \land \ldots \land C_{m} \\
	\text{ Where } C_{i} &= ( l_1 \lor l_2 \lor \ldots \lor l_k)\\
	\text{ Where } l_{j} &= p \text{ or } l_{j} = \lnot p
\end{align*}
Each $C_i$ is called a clause, and each $l_j$ is called a literal. One can observe that $C$ does not belong to FORM. Let us consider an example to understand the scenario:
let $\varphi := (C_1 \land (C_2 \land C_3))$ and $\psi := ((C_1 \land C_2) \land C_3)$. As models($\varphi$) = models($\psi$), $\varphi \equiv \psi$. Therefore, as it does not matter where we put the parenthesis in $C$, we just omit the use of it.

Let us find the models of formula $C$.
\begin{enumerate}
	\item models($C$) = $\bigcap$ models($C_i$).
	\item models($C_i$) = $\bigcup$ models($l_j$).
	\item If $C = \{\}$, then models($C$) = $2^{\text{PROP}}$.
	\item If $C_{i} = \{\}$, then models($C_i$) = $\emptyset$. Let us denote an empty clause as $\square$, that is models($\square$) = $\emptyset$.
	\item if $C = \{ \square\}$, then  models($C$) = $\emptyset$.

\end{enumerate}

\textit{Can we express arbitrary formula in CNF?}
\begin{theorem}
Every formula $\varphi$ can be represented in $\varphi_{\text{CNF}}$.
\end{theorem}

We can apply the distribution law to convert $\varphi$ to $\varphi_{\text{CNF}}$. But, in the worst case, it may take exponential many steps.

Here is an example that uses distribution law to compute the CNF representation $\varphi_{\text{CNF}}$ of a formula $\varphi$:

\begin{align*}
\varphi &= ((y_1 \land y_2) \lor (y_3 \land y_4)) \lor (y_5 \land y_6)\\
		&\equiv ((y_1 \lor (y_3 \land y_4))  \land (y_2 \lor (y_3 \land y_4))) \lor (y_5 \land y_6)\\
		&\equiv (((y_1 \lor y_3) \land (y_1 \lor y_4))  \land ((y_2 \lor y_3) \land (y_2 \lor y_4))) \lor (y_5 \land y_6)\\
		&\equiv (y_1 \lor y_3 \lor y_5)  \land (y_1 \lor y_3 \lor y_6) \land (y_1 \lor y_4 \lor y_5)  \land (y_1 \lor y_4 \lor y_6) \land (y_2 \lor y_3 \lor y_5)  \\
		& \quad \quad \land (y_2 \lor y_3 \lor y_6) \land (y_2 \lor y_4 \lor y_5)  \land (y_2 \lor y_4 \lor y_6) \\
		&\equiv \varphi_{\text{CNF}}
\end{align*}


Now, the question is, \emph{given a formula $\varphi$, can we determine whether it is satisfiable }

 \begin{algorithm}[H]

 	\caption{CheckSAT{($\varphi$)}\label{algo:checksat}}
 	\begin{algorithmic}[1]
 		\For { $\tau \in 2^{\text{PROP}}$}\label{line:checksat:prop}
	 		\If {$\varphi(\tau) = 1 $}
		 		\State
		 		\Return {{SAT, $\tau$}}

	 		\EndIf
 		\EndFor
 		\State
 		\Return {{UNSAT, $\emptyset$}}
 	\end{algorithmic}
 \end{algorithm}
Algorithm~\ref{algo:checksat} iterates over $2^{\text{PROP}}$ to find a $\tau$, such that $\varphi(\tau)=1$. We can do slightly better by iterating over $\ap{(\varphi)}$. But, again $\ap{(\varphi)}$ can be PROP. In the worst case scenario, finding a such assignment of $\varphi$ can take exponential time with Algorithm~\ref{algo:checksat}.

\emph{Can we do better?} We don't know.
This is the fundamental question of our time, and this question has been at the heart of computer science ever since computer science has been a field.
We do not know, for any given formula $\varphi$, does there exist an efficient (i.e. polynomial-time) algorithm which always finds an assignment
over $\ap{(\varphi)}$ that evaluates $\varphi$ to be true, whenever it exists.

However, we can still look for methods that run in a reasonable amount of time for most formulas, even though for some worst-case formulas they might take exponential amounts of time.

\subsection{Resolution Refutation}
\begin{definition}
Resolution:  $(p \lor \alpha)\land (\lnot p \lor \beta) \iff (\alpha \lor \beta)$
\end{definition}
Let $\varphi = (p \lor \alpha)\land (\lnot p \lor \beta)$, and $\psi = (\alpha \lor \beta)$. Observe that $\varphi$ and $\psi$ are not equivalent, however $\varphi$ is satisfiable if and only if $\psi$ is satisfiable, formulas $\varphi$ and $\psi$ are called equisatisfiable formulas. The real benefit of resolution is if we somehow know that $(\alpha \lor \beta)$ is unsatisfiable then $(p \lor \alpha)\land (\lnot p \lor \beta)$ is also unsatisfiable.

Now that, we are only interested in equisatisfiable formulas, let us again discuss the representation in CNF.

If we are only interested in satisfiability of $\varphi$, then we can convert it to an equisatisfiable formula $\varphi_{\text{CNF}}$ in polynomial time.

\begin{align*}
\varphi &= ((y_1 \land y_2) \lor (y_3 \land y_4)) \lor (y_5 \land y_6)\\
&\equiv (t_1 \implies (y_1 \land y_2)) \land (t_2 \implies (y_3 \land y_4)) \land (t_3 \implies (y_5 \land y_6)) \land (t_1 \lor t_2 \lor t_3)\\
&\equiv (\lnot t_1 \lor y_1) \land (\lnot t_1 \lor y_2) \land (\lnot t_2 \lor y_3) \land (\lnot t_2 \lor y_4) \land (\lnot t_3 \lor y_5 ) \land (\lnot t_3 \lor y_6) \land (t_1 \lor t_2 \lor t_3)\\
&\equiv \varphi_{\text{CNF}}
\end{align*}

The above transformation was proposed by Tseytin, and called as Tseytin transformation.
\begin{theorem}
	Every formula $\varphi$ can be represented to $\varphi_{\text{CNF}}$ in polynomial time such that $\varphi$ is satisfiable if and only if $\varphi_{\text{CNF}}$ is satisfiable.
\end{theorem}

Hence, we can determine whether $\varphi_{\text{CNF}}$ is satisfiable, then we know whether $\varphi$ is satisfiable. We are interested in CNF formula, as it directly give a method based on resolution to determine the satisfiability of formula.

\begin{definition}
	List of clauses $C_1, C_2,\ldots C_{t}$ is a \textbf{resolution refutation} of formula $\varphi_{\text{CNF}}$ if:
	\begin{enumerate}
		\item $C_t := \square$.
		\item $C_k \in \varphi_{\text{CNF}}$ or $C_k$ is derived using resolution from $C_i, C_j$, where $i,j < k$.
	\end{enumerate}
\end{definition}

\subsubsection{Power of Resolution Refutation}
\begin{theorem}
	\label{theorem:resolution_proof}
	A formula $\varphi_{\text{CNF}}$ is refutable if and only if $\varphi_{\text{CNF}}$ is unsatisfiable.
\end{theorem}
\end{document}
